% The spoken body of the homily for the Nineteenth Sunday after Pentecost.
% Continuous preaching, ready to read aloud: no heading, no direction to the
% preacher, no citation inside the speech. The loci, the relation to the
% reviewed readings, the audience, the word count and the pace estimate stand
% in the terminal note, which no one reads out. The blank lines are the
% paragraph breaks a speaker recovers his place by, and the seven vertical
% spaces divide the speech into its eight movements: the guest who was
% silent; how the hall was filled, with Chrysostom and the Collect; what the
% garment is not, and what it is, with Gregory and Augustine; the Epistle's new
% man with Jerome, its four acts, the bound hands and the Gradual's lifted
% hands with Hilary; the evening sacrifice of the Cross with Augustine, and the
% altar; the garment asked for, in the Communion with Augustine, Hilary's
% baptismal gift, the Introit and the Postcommunion; the two feasts, the two
% wishes and the week's response; and the hope of the end. The columns are set
% ragged at the foot, so that the seven breaks keep one height instead of
% absorbing the columns' balancing.
\raggedcolumns

``Friend, how camest thou in hither not having on a wedding garment?''

The king asks that question near the end of today's Gospel. He asks it of a man
who had done almost everything right. He was invited, and he came. He did not
go off to his farm or his business, or lay hands on the servants. He sat down
at the king's table. Then the king came in to see his guests, called this man
``Friend'', and asked him one question. ``But he was silent.''

Of everyone in the parable, that man is the most like us. We were invited, and
we have come. We are inside. So his question is ours: what should he have been
wearing?

\bigskip

Remember how the hall was filled. The first guests ``would not come''. Called
again, they ``went their ways, one to his farm and another to his
merchandise'', and the rest killed the servants. So the king sent his servants
into the highways, and they gathered ``both bad and good: and the marriage was
filled with guests.''

St~John Chrysostom grants that the excuses seem reasonable, and draws the
lesson: ``though the things which hinder us be necessary, to set the things
spiritual at a higher price than all.'' Today's Collect asks for the
freedom to do it: that ``being free both in soul and body'' we may ``perform
thy service''. And Chrysostom reminded his people: ``Hear whence ye were
called. From the highway.'' We were found on the road, and brought in.

\bigskip

What, then, is the wedding garment? St~Gregory the Great asked the people of
Rome: if it is baptism, or faith, who ever came into this wedding without them?
Whatever the man lacks, it is not what brought him in.

St~Augustine asked whether it might be the altar, or what is received there,
and answered: ``But no; we see that many eat, and `eat and drink judgment to
themselves.'\,'' The garment, he says, ``is in the heart, not on the body''.
The servants never noticed the man. The king did.

Both give one answer: the garment is charity. Gregory gives the reason. Charity
is rightly called the wedding garment, he says, because our Maker himself had
it when he came to the wedding at which he joined the Church to himself. The
garment is the Bridegroom's own dress. The king calls the man ``Friend'' and
still condemns him, and Gregory hears in that word a friend by faith, and no
friend by what he has done.

\bigskip

Listen again to the Epistle. ``Put on the new man,'' says St~Paul, ``who
according to God is created in justice and holiness of truth.'' The Epistle
commands us to be clothed; the Gospel ends with a man who is not. St~Jerome
joins the two himself. The wedding garment, he writes, is the Lord's
commandments and the works that make up the garment of the new man. And on this
same Epistle he says who the new man is: Christ, with whom all believers ought
to be clothed.

St~Paul gives the new man four plain acts. ``Putting away lying, speak ye the
truth, every man with his neighbour. For we are members one of another.'' ``Be
angry: and sin not. Let not the sun go down upon your anger.'' ``Give not place
to the devil.'' And the thief: ``let him labour, working with his hands the
thing which is good, that he may have something to give to him that suffereth
need.''

Truth to your neighbour. Anger ended before sunset. No foothold for the devil.
Hands that work in order to give. That is what the wedding garment looks like
on an ordinary weekday.

Look at the hands. The king has the man bound hand and foot. St~Gregory says
the binding had begun long before. Hands that gave nothing to the needy, and
feet that would not visit the sick, were already bound from doing good, by
their own will. Between the lessons, the Gradual prays for ``the lifting up of
my hands, as evening sacrifice.'' St~Hilary of Poitiers says that to spread the
hands is the posture of prayer, but to lift them is to do the work. The evening
sacrifice that pleases God, he says, is the one our Lord names at the judgment:
``I was hungry, and you gave me to eat.''

\bigskip

St~Augustine hears the Gradual's verse differently: first of Christ, who laid down his
life upon the Cross as the day sank towards evening. He says: ``That then is
the `evening sacrifice,' the Passion of the Lord, the Cross of the Lord''; and
it ``produced, in His Resurrection, a morning offering.'' Hilary looks at our
hands, Augustine at the Lord's, and we need both. The charity the king looks
for in us is the charity his Son wore on the Cross, ``praying to His Father for
His enemies'', as Augustine says.

That sacrifice is not only remembered here. On this altar the one sacrifice of
the Cross is made present, and the risen Christ gives his Body and Blood to
those who come to his table.

\bigskip

But a fair question rises. If the guests were swept in from the highways, how
could any of them have had a garment? And if the garment is charity, which of
us can weave it alone?

At the Communion we will hear: ``Thou hast commanded thy commandments to be
kept most diligently. O! that my ways may be directed to keep thy
justifications.'' A command, and then a wish. Of that wish St~Augustine says:
``When thou hearest, `O that,' recognise the words of one wishing; and having
recognised the expression of a wish, lay aside the pride of presumption.'' And
he concludes: ``God must be prayed to grant Himself what He enjoineth.''

St~Hilary looks at it from the other side. For him the garment is the glory of
the Holy Spirit, received at baptism, to be kept spotless until the kingdom. He
and Gregory differ over what to call it and where it begins; but both mean
something the baptized can lose, or never put on.

So the garment is God's gift, and we are still to put it on and keep it: by
doing, and by asking. Augustine says both: ``Clothe others, and ye are clothed
yourselves''; and, ``Run to Him, beseech Him; He knoweth how to sanctify His
faithful ones, He knoweth how to clothe His naked ones.'' The Mass opened with
the Lord's promise: ``in whatever distress they call on me, I will hear them''.
And the prayer after Communion asks that God's healing work ``make us always
obedient to thy commandments''.

\bigskip

Augustine taught that there are ``two feasts of the Lord; one to which the good
and evil come, the other to which the evil come not.'' The feast of today's
Gospel, he said, is the first; he called it the Lord's Table and the Feast of
the Holy Scriptures. Of it he said: ``Be that feast which now is, received
worthily, that we may attain to the other.'' And he warned those coming to the
Lord's Table not to pray: ``O that mine enemy might die!''

So there are two wishes, each beginning ``O that''. One asks God to destroy
whoever has hurt us. The other asks God to set our own ways straight. Which
will we bring to the altar today?

Before you come, ask yourself St~Gregory's questions: whether you hold hatred
against anyone; whether envy of someone else's happiness burns in you. Then,
this week, do three things. If you are angry with someone, settle it before the
sun goes down, or at least forgive him before God today. Let your hands give
something that costs you: food, clothing, an hour with someone sick. And every
morning, pray the Communion's words: ``O! that my ways may be directed to keep
thy justifications.''

\bigskip

The king in the parable came in once and found one guest unready. He has not
yet come in for us. St~Gregory says that we know we are called, but not whether
we are chosen. That is a reason to be humble, not to despair. For to anyone who
wears the garment imperfectly, Gregory offers the psalm's hope: ``Thy eyes did
see my imperfect being.''

So let us come to the table clothed in charity: the new man, Christ himself put
on. When the king comes in to see his guests, may he find his Son's own dress
on us. And may he say to us, without reproach, the word he said to that silent
man: ``Friend''. The one who calls us is the one who clothes us. He knows how
to clothe his naked ones.
